
\documentclass{article}
\usepackage{colortbl}
\usepackage{makecell}
\usepackage{multirow}
\usepackage{supertabular}

\begin{document}

\newcounter{utterance}

\centering \large Interaction Transcript for game `codenames', experiment `association\_easy', episode 0 with Qwen3.6{-}35B{-}A3B{-}FP8{-}without{-}reasoning.
\vspace{24pt}

{ \footnotesize  \setcounter{utterance}{1}
\setlength{\tabcolsep}{0pt}
\begin{supertabular}{c@{$\;$}|p{.15\linewidth}@{}p{.15\linewidth}p{.15\linewidth}p{.15\linewidth}p{.15\linewidth}p{.15\linewidth}}
    \# & $\;$A & \multicolumn{4}{c}{Game Master} & $\;\:$B\\
    \hline

    \theutterance \stepcounter{utterance}  
    & & & \multicolumn{2}{p{0.3\linewidth}}{
        \cellcolor[rgb]{0.9,0.9,0.9}{
            \makecell[{{p{\linewidth}}}]{
                \texttt{\tiny{[GM$|$GM]}}
                {'hidden': {'team': ['sarkans', 'zils', 'āmurs', 'priesteris', 'īlens', 'iedegums', 'urbis', 'imams', 'sludinātājs'], 'innocent': ['vara galva', 'melns', 'volejbols', 'ciklons', 'žurka', 'cīņas sports', 'mākoņi'], 'opponent': ['putotājs', 'mētelis', 'standziņas', 'šķīvis', 'džemperis', 'krekls', 'sietiņš', 't-krekls'], 'assassin': ['apšuvums']}, 'revealed': {'team': {'team': [], 'innocent': [], 'opponent': [], 'assassin': []}, 'opponent': {'team': [], 'innocent': [], 'opponent': [], 'assassin': []}}}
            }
        }
    }
    & & \\ \\

    \theutterance \stepcounter{utterance}  
    & & \multicolumn{4}{p{0.6\linewidth}}{
        \cellcolor[rgb]{0.9,0.9,0.9}{
            \makecell[{{p{\linewidth}}}]{
                \texttt{\tiny{[P1$\langle$GM]}}
                \texttt{Sniedziet vienu vienvārdu norādi, kas ir saistīta ar kādu no jūsu komandas vārdiem no tālāk dotā saraksta. Jūs varat izvēlēties, uz cik un kādiem vārdiem attiecas jūsu norāde.} \\
\texttt{Mērķis ir panākt, lai jūsu partneris uzmin visus jūsu komandas vārdus, pamatojoties uz jūsu dotajām norādēm, pirms otra komanda ir atklājusi visus savus vārdus. Ja jūsu partneris uzmin "pretinieka" vai "novirzītāja" vārdu, viņa uzminēšanas kārta beidzas, un pretinieku komanda var veikt savu kārtu. Ja jūsu partneris uzmin "nogalināšanas" vārdu, spēle beidzas un jūsu komanda uzreiz zaudē.} \\
\texttt{Jo vairāk vārdu, uz kuriem attiecas jūsu norāde, jo ātrāk jūs varat uzvarēt spēli, bet jo grūtāk asociācijas var būt uzminēt jūsu partnerim, tāpēc izvēlieties savu norādes vārdu gudri. Norādes vārdam ir jābūt semantiski saistītam ar mērķa vārdiem, tas nedrīkst būt viens no sarakstos iekļautajiem vārdiem vai saturēt to daļas.} \\
\\ 
\texttt{Vienvārdu norādi un ar komatu atdalīto saistīto mērķa vārdu sarakstu vienmēr sniedziet šādā formātā un atbildi veidojiet pēc iespējas īsāku, nekad neiekļaujiet citu tekstu, kas nav prasīts šajā formā:} \\
\\ 
\texttt{PADOMS: <VĀRDS>} \\
\texttt{MĒRĶI: <MĒRĶI>} \\
\\ 
\texttt{Jūsu komandas vārdi ir: sarkans, zils, āmurs, priesteris, īlens, iedegums, urbis, imams, sludinātājs.} \\
\texttt{Jūsu pretinieka komandas vārdi ir: putotājs, mētelis, standziņas, šķīvis, džemperis, krekls, sietiņš, t{-}krekls.} \\
\texttt{Novirzītāja vārdi ir: vara galva, melns, volejbols, ciklons, žurka, cīņas sports, mākoņi.} \\
\texttt{Nogalināšanas vārdi ir: apšuvums.} \\
            }
        }
    }
    & \\ \\

    \theutterance \stepcounter{utterance}  
    & \multicolumn{4}{p{0.6\linewidth}}{
        \cellcolor[rgb]{0.9,0.9,0.9}{
            \makecell[{{p{\linewidth}}}]{
                \texttt{\tiny{[P1$\rangle$GM]}}
                \texttt{PADOMS: krāsa} \\
\texttt{MĒRĶI: sarkans, zils} \\
            }
        }
    }
    & & \\ \\

    \theutterance \stepcounter{utterance}  
    & & & \multicolumn{2}{p{0.3\linewidth}}{
        \cellcolor[rgb]{0.9,0.9,0.9}{
            \makecell[{{p{\linewidth}}}]{
                \texttt{\tiny{[GM$|$GM]}}
                \texttt{krāsa} \\
            }
        }
    }
    & & \\ \\

    \theutterance \stepcounter{utterance}  
    & & & \multicolumn{2}{p{0.3\linewidth}}{
        \cellcolor[rgb]{0.9,0.9,0.9}{
            \makecell[{{p{\linewidth}}}]{
                \texttt{\tiny{[GM$|$GM]}}
                ['sarkans', 'zils']
            }
        }
    }
    & & \\ \\

    \theutterance \stepcounter{utterance}  
    & & & \multicolumn{2}{p{0.3\linewidth}}{
        \cellcolor[rgb]{0.9,0.9,0.9}{
            \makecell[{{p{\linewidth}}}]{
                \texttt{\tiny{[GM$|$GM]}}
                {'word': 'sarkans', 'assignment': 'team'}
            }
        }
    }
    & & \\ \\

    \theutterance \stepcounter{utterance}  
    & & & \multicolumn{2}{p{0.3\linewidth}}{
        \cellcolor[rgb]{0.9,0.9,0.9}{
            \makecell[{{p{\linewidth}}}]{
                \texttt{\tiny{[GM$|$GM]}}
                {'word': 'zils', 'assignment': 'team'}
            }
        }
    }
    & & \\ \\

    \theutterance \stepcounter{utterance}  
    & & \multicolumn{4}{p{0.6\linewidth}}{
        \cellcolor[rgb]{0.9,0.9,0.9}{
            \makecell[{{p{\linewidth}}}]{
                \texttt{\tiny{[GM$\rangle$P2]}}
                \texttt{Sniedziet ar komatiem atdalītu sarakstu ar līdz 2 vārdiem no turpmāk minētā saraksta, kas vislabāk attiecas uz vārdu 'krāsa' vai ir visciešāk saistīti ar to. Vienmēr sāciet minējumu sarakstu ar MINĒJUMS:  un atbildē neiekļaujiet nekādu citu tekstu.} \\
\\ 
\texttt{t{-}krekls, šķīvis, džemperis, imams, vara galva, iedegums, urbis, žurka, krekls, volejbols, standziņas, mākoņi, sludinātājs, cīņas sports, melns, sarkans, sietiņš, mētelis, zils, āmurs, ciklons, priesteris, putotājs, apšuvums, īlens} \\
            }
        }
    }
    & \\ \\

    \theutterance \stepcounter{utterance}  
    & & & \multicolumn{4}{p{0.6\linewidth}}{
        \cellcolor[rgb]{0.9,0.9,0.9}{
            \makecell[{{p{\linewidth}}}]{
                \texttt{\tiny{[GM$\langle$P2]}}
                \texttt{MINĒJUMS: melns, sarkans, zils, iedegums, vara galva} \\
            }
        }
    }
     \\ \\

    \theutterance \stepcounter{utterance}  
    & & & \multicolumn{2}{p{0.3\linewidth}}{
        \cellcolor[rgb]{0.9,0.9,0.9}{
            \makecell[{{p{\linewidth}}}]{
                \texttt{\tiny{[GM$|$GM]}}
                {'player': 'Guesser', 'type': 'wrong number of guesses', 'utterance': 'melns, sarkans, zils, iedegums, vara galva', 'message': 'You made too many guesses (5). You are only allowed to make 2 guesses!', 'guesses': ['melns', 'sarkans', 'zils', 'iedegums', 'vara galva'], 'number_of_allowed_guesses': 2}
            }
        }
    }
    & & \\ \\

\end{supertabular}
}

\end{document}
